\documentclass[12pt,a4paper,final]{article}

\usepackage[T2A]{fontenc}
\usepackage[utf8]{inputenc}
\usepackage[english, russian]{babel}
\usepackage{amsmath} 
\usepackage{fancyhdr}
\usepackage[usenames]{color}
\usepackage{colortbl}
\usepackage[left=3.00cm, right=2.00cm,top=3.00cm,bottom=3.00cm]{geometry}
\newcommand{\ds}{\displaystyle}
\usepackage{graphicx}
\graphicspath{{img/}{plots/}}
\usepackage{tikz}
\usepackage{setspace}
\onehalfspacing
\usepackage{float}
\usepackage{appendix}
\usepackage{amssymb}
\usepackage{pdfpages}
\usepackage{listings}
\usepackage{moreverb}
\usepackage{pdflscape}  % Лучше, чем lscape, для PDF
\usepackage{booktabs}
\usepackage{longtable}
\newcommand{\req}[2]{\par\smallskip\noindent\textbf{#1}\hspace{0.6em}#2\par}

% ---------- Пакеты и команды для таблиц и списков требований ----------
\usepackage{array}      % типы столбцов C{..} и L{..}
\usepackage{longtable}  % многостраничные таблицы
\usepackage{enumitem}   % настройка списков (окружение reqs)
\usepackage{caption}
\captionsetup{labelsep=endash,justification=raggedright,singlelinecheck=false}
\captionsetup[longtable]{labelsep=endash,justification=raggedright,singlelinecheck=false}
\renewcommand{\arraystretch}{1.15}
\newcolumntype{C}[1]{>{\centering\arraybackslash}p{#1}}   % по центру, ширина #1
\newcolumntype{L}[1]{>{\raggedright\arraybackslash}p{#1}} % по левому краю, ширина #1

% Список требований с идентификаторами (ОБЩ-01, ФЛТ-02, ...)
\newenvironment{reqs}
  {\begin{description}[style=sameline,leftmargin=2.2cm,labelwidth=2.0cm,
     labelsep=0.2cm,font=\normalfont\bfseries,itemsep=3pt]}
  {\end{description}}
\newcommand{\src}[1]{\textit{(п.~#1 задания)}}
\newcommand{\der}{\textit{(производное требование)}}
\providecommand{\textenglish}[1]{\foreignlanguage{english}{#1}}

% Заголовок приложения: #1 — буква, #2 — название
\newcommand{\appsection}[2]{%
  \clearpage
  \phantomsection
  \addcontentsline{toc}{section}{Приложение #1. #2}%
  \begin{center}
    {\large\bfseries ПРИЛОЖЕНИЕ #1}\\
    (обязательное)\\[2pt]
    {\large\bfseries #2}
  \end{center}}
% ----------------------------------------------------------------------

\usepackage[unicode, pdftex, hidelinks]{hyperref}

\definecolor{apricot}{HTML}{FFF0DA}
\definecolor{mygreen}{rgb}{0,0.6,0}
\definecolor{string}{HTML}{B40000} % цвет строк в коде
\definecolor{comment}{HTML}{008000} % цвет комментариев в коде
\definecolor{keyword}{HTML}{1A00FF} % цвет ключевых слов в коде
\definecolor{morecomment}{HTML}{8000FF} % цвет include и других элементов в коде
\definecolor{captiontext}{HTML}{FFFFFF} % цвет текста заголовка в коде
\definecolor{captionbk}{HTML}{999999} % цвет фона заголовка в коде
\definecolor{bk}{HTML}{FFFFFF} % цвет фона в коде
\definecolor{frame}{HTML}{999999} % цвет рамки в коде
\definecolor{brackets}{HTML}{B40000} % цвет скобок в коде


\usepackage{minted}
% подгружаемые языки — подробнее в документации listings
\lstloadlanguages{Matlab}
\lstset{
  language=Matlab, % Язык кода по умолчанию
  morekeywords={*,...}, % если хотите добавить ключевые слова, то добавляйте
  % Цвета
  keywordstyle=\color{keyword}\ttfamily\bfseries,
  %stringstyle=\color{string}\ttfamily,
  stringstyle=\ttfamily\color{red!50!brown},
  commentstyle=\color{comment}\ttfamily,
  morecomment=[l][\color{morecomment}]{\#},
  % Настройки отображения
  breaklines=true, % Перенос длинных строк
  basicstyle=\ttfamily\footnotesize, % Шрифт для отображения кода
  backgroundcolor=\color{bk}, % Цвет фона кода
  frame=lrb,xleftmargin=\fboxsep,xrightmargin=-\fboxsep, % Рамка, подогнанная к заголовку
  rulecolor=\color{frame}, % Цвет рамки
  tabsize=3, % Размер табуляции в пробелах
  % Настройка отображения номеров строк. Если не нужно, то удалите весь блок
  numbers=left, % Слева отображаются номера строк
  stepnumber=1, % Каждую строку нумеровать
  numbersep=5pt, % Отступ от кода
  numberstyle=\small\color{black}, % Стиль написания номеров строк
  % Для отображения русского языка
  extendedchars=true,
  literate={Ö}{ {\'O} }1
  {~}{ {\textasciitilde} }1
  {а}{ {\selectfont\char224} }1
  {б}{ {\selectfont\char225} }1
  {в}{ {\selectfont\char226} }1
  {г}{ {\selectfont\char227} }1
  {д}{ {\selectfont\char228} }1
  {е}{ {\selectfont\char229} }1
  {ё}{ {\'e} }1
  {ж}{ {\selectfont\char230} }1
  {з}{ {\selectfont\char231} }1
  {и}{ {\selectfont\char232} }1
  {й}{ {\selectfont\char233} }1
  {к}{ {\selectfont\char234} }1
  {л}{ {\selectfont\char235} }1
  {м}{ {\selectfont\char236} }1
  {н}{ {\selectfont\char237} }1
  {о}{ {\selectfont\char238} }1
  {п}{ {\selectfont\char239} }1
  {р}{ {\selectfont\char240} }1
  {с}{ {\selectfont\char241} }1
  {т}{ {\selectfont\char242} }1
  {у}{ {\selectfont\char243} }1
  {ф}{ {\selectfont\char244} }1
  {х}{ {\selectfont\char245} }1
  {ц}{ {\selectfont\char246} }1
  {ч}{ {\selectfont\char247} }1
  {ш}{ {\selectfont\char248} }1
  {щ}{ {\selectfont\char249} }1
  {ъ}{ {\selectfont\char250} }1
  {ы}{ {\selectfont\char251} }1
  {ь}{ {\selectfont\char252} }1
  {э}{ {\selectfont\char253} }1
  {ю}{ {\selectfont\char254} }1
  {я}{ {\selectfont\char255} }1
  {А}{ {\selectfont\char192} }1
  {Б}{ {\selectfont\char193} }1
  {В}{ {\selectfont\char194} }1
  {Г}{ {\selectfont\char195} }1
  {Д}{ {\selectfont\char196} }1
  {Е}{ {\selectfont\char197} }1
  {Ё}{ {\'E} }1
  {Ж}{ {\selectfont\char198} }1
  {З}{ {\selectfont\char199} }1
  {И}{ {\selectfont\char200} }1
  {Й}{ {\selectfont\char201} }1
  {К}{ {\selectfont\char202} }1
  {Л}{ {\selectfont\char203} }1
  {М}{ {\selectfont\char204} }1
  {Н}{ {\selectfont\char205} }1
  {О}{ {\selectfont\char206} }1
  {П}{ {\selectfont\char207} }1
  {Р}{ {\selectfont\char208} }1
  {С}{ {\selectfont\char209} }1
  {Т}{ {\selectfont\char210} }1
  {У}{ {\selectfont\char211} }1
  {Ф}{ {\selectfont\char212} }1
  {Х}{ {\selectfont\char213} }1
  {Ц}{ {\selectfont\char214} }1
  {Ч}{ {\selectfont\char215} }1
  {Ш}{ {\selectfont\char216} }1
  {Щ}{ {\selectfont\char217} }1
  {Ъ}{ {\selectfont\char218} }1
  {Ы}{ {\selectfont\char219} }1
  {Ь}{ {\selectfont\char220} }1
  {Э}{ {\selectfont\char221} }1
  {Ю}{ {\selectfont\char222} }1
  {Я}{ {\selectfont\char223} }1
  {\{}{ { {\color{brackets}\{} } }1 % Цвет скобок {
  {\} }{ { {\color{brackets}\} } } }1 % Цвет скобок }
}
\begin{document}
	\noindent
	\begin{center}
		\thispagestyle{empty}
		\large{МИНИСТЕРСТВО НАУКИ И ВЫСШЕГО ОБРАЗОВАНИЯ РОССИИ\\
			\vspace{3mm}
			ФЕДЕРАЛЬНОЕ ГОСУДАРСТВЕННОЕ БЮДЖЕТНОЕ ОБРАЗОВАТЕЛЬНОЕ}\end{center}\vspace{-5mm}
	\begin{center}\large{УЧРЕЖДЕНИЕ}\end{center}
	\begin{center}\vspace{-3mm}\large{ВЫСШЕГО ПРОФЕССИОНАЛЬНОГО ОБРАЗОВАНИЯ}\end{center}
	\begin{center}\large\textbf{<<САНКТ-ПЕТЕРБУРГСКИЙ ПОЛИТЕХНИЧЕСКИЙ УНИВЕРСИТЕТ
			ПЕТРА ВЕЛИКОГО>>}\end{center}
	\begin{center}
		\normalsize{Институт Компьютерных наук и Кибербезопасности}
	\end{center}
	\begin{center}
		Высшая школа технологий искусственного интеллекта
	\end{center}
	\begin{center}
		Направление 02.03.01 Математика и компьютерные науки
	\end{center}
	
	\vspace{0.3cm}
	
	\begin{center}
		\large Защита информации
	\end{center}
	\begin{center}
		\Large\textbf{Спецификация требований} 
	\end{center}
	\begin{center}
		\large<<Разработка межсетевого экрана>> 
        \end{center}

	
	\begin{flushleft}
		Студенты группы
	\end{flushleft}
	\vspace{-0.55cm}
	\begin{flushleft}
		\hbox to 5.5cm {№5130201/30102:\quad \hrulefill}
	\end{flushleft}
	\begin{flushright}
		\vspace{-1.45cm}
		Смирнов С.Р.
	\end{flushright}
    \begin{flushright}
		
		Михайлова А.А.
	\end{flushright}
    \begin{flushright}
		
		Аббад М.С.
	\end{flushright}
	\vspace{1cm}
	\begin{flushleft}
		\hbox to 5.5cm {Преподаватель:\quad \hrulefill}
	\end{flushleft}
	\begin{flushright}
		\vspace{-1.5cm}
		Силиненко А.В.
	\end{flushright}
	\begin{flushright}
		\vspace{1cm}
		\hspace{5cm}\hbox to 1.2cm {<<\hrulefill>>}\hbox to 2.2cm {\hrulefill }\hbox to 1.5cm { 20\hrulefill г.}
	\end{flushright}
	\begin{center}
		\vspace{4cm}
		Санкт-Петербург 2026
	\end{center}
	\newpage
    \section*{Состав исполнителей и распределение ответственности}
 
\noindent Группа 5130201/30102. Состав команд приведён по уточнённому списку
участников проекта. Каждый исполнитель участвует в реализации и проверке
требований закреплённой команды.
 
{\singlespacing
\renewcommand{\arraystretch}{1.1}
\noindent
\begin{tabular}{@{}L{2.9cm}L{3.6cm}L{9.3cm}@{}}
\toprule
Команда / подсистема & Исполнители & Область ответственности\\
\midrule
Фильтрация & Алехин М.\,И.\newline Романова А.\,П.\newline Сафин А.\,Р. &
Приём и передача кадров, политика доступа, таблица соединений;
ФЛТ-01--10; СОС-01--06; ОБЩ-02, ОБЩ-03; НДЖ-02; ОБР-01, ОБР-03--ОБР-13.\\
\addlinespace[6pt]
Регистрация &
Невечерин А.\,А.\newline Саржан А.\,А.\newline Гладков И.\,А.\newline Коновалов Д.\,С.  &
Журналы событий и пакетов, RAM-буферы, сохранение и ротация;
РЕГ-01--07; УПР-04--07.\\
\addlinespace[6pt]
Управление &
Гвоздева Е.\,В.\newline Козловская О.\,А.\newline Фролов К.\,А. \newline Крутиков Д.\,Д. &
Независимые CLI и web, отдельная авторизация, учётные записи и сеансы;
ДСТ-01--08; УПР-01, УПР-02, УПР-12; интерфейс ко всем командам приложения~Б.\\
\addlinespace[6pt]
Контроль\newline целостности &
Путята М.\,А.\newline Смирнов С.\,Р.\newline Аббад М.\,С.\newline Михайлова А.\,А. &
ЦЛС-01--07; системный модуль, настройка ОС (ОБР-02); координация конфигурации; УПР-09--11,
УПР-13--16; НДЖ-01, НДЖ-03--05; сборка, установка и испытания
в аудитории~307.\\
\bottomrule
\end{tabular}
\par}
 
\medskip
Общие требования ОБЩ-01 и совместимость API обеспечиваются всеми
командами. Аналитическая работа по границам подсистем выполняется совместно:
согласуются структуры данных, команды, ошибки и сценарии взаимодействия.
Владелец подсистемы отвечает за реализацию своего API, а команды-потребители~—
за совместные проверки.
 
Реализация web и CLI ведётся раздельно, хотя исполнители входят в одну команду
управления. Подсистема авторизации также самостоятельна; она не является
отдельной копией базы в каждом интерфейсе.
\clearpage
\tableofcontents
\clearpage

% ============================== 1 ВВЕДЕНИЕ ============================
\section{Введение}

\subsection{Наименование программы}

Полное наименование: «Программное обеспечение межсетевого экрана».
Краткое наименование: ПО МЭ.

\subsection{Краткая характеристика области применения}

ПО МЭ предназначено для установки на вычислительное устройство под управлением
ОС Linux, имеющее от одного до шести портов Ethernet, и обеспечивает защиту
информации в сегментах локальных вычислительных сетей, подключённых к портам
устройства, путём контроля и фильтрации проходящего через устройство сетевого
трафика на канальном, сетевом и транспортном уровнях в соответствии с политикой
доступа, задаваемой администратором.

ПО МЭ может применяться для разграничения доступа между сегментами сети
организации, для защиты внутренних сегментов от доступа из внешних сетей,
а также для трансляции сетевых адресов на границе сегментов.

\subsection{Термины, определения и сокращения}

В настоящем документе применяются следующие термины и сокращения.

\begin{description}
\item[МЭ (межсетевой экран)] — программное или программно-техническое средство,
  реализующее контроль и фильтрацию сетевого трафика в соответствии с заданными
  правилами.
\item[Изделие] — вычислительное устройство с установленным ПО МЭ.
\item[Интерфейс (порт)] — сетевой интерфейс Ethernet изделия, обслуживаемый ПО МЭ.
\item[Политика доступа] — упорядоченная совокупность правил фильтрации,
  определяющая порядок обработки пакетов.
\item[Правило фильтрации (правило)] — элемент политики доступа, содержащий
  критерии отбора пакетов и атрибуты, определяющие обработку отобранных пакетов.
\item[Контроль состояния] — механизм отслеживания сеансов взаимодействия
  (соединений) и обработки пакетов с учётом принадлежности их к установленным
  сеансам (англ. \textenglish{stateful inspection}).
\item[Таблица состояний] — совокупность записей о сеансах, отслеживаемых
  механизмом контроля состояния.

\item[NAT] — трансляция сетевых адресов (англ. Network Address Translation); SNAT —
трансляция адреса источника; DNAT — трансляция адреса приёмника.
\item[Правило NAT] - элемент таблицы DNAT или SNAT, содержащий критерии отбора
пакетов и параметры трансляции адреса и (или) порта. Правило NAT не
определяет, будет ли пакет передан или удалён.
 
\item[Таблица DNAT, таблица SNAT] - независимые упорядоченные списки правил
NAT, применяемых соответственно к адресу и порту приёмника и к адресу и
порту источника.
 
\item[Преобразование портов (PAT)] - замена номера TCP/UDP-порта источника или
приёмника, выполняемая правилом NAT совместно с заменой IPv4-адреса либо
отдельно от неё.
 
\item[API управления] - программный интерфейс подсистемы управления, через
который командный и web-интерфейсы читают и изменяют политику доступа и
правила NAT.
\item[Администратор] — пользователь, уполномоченный на управление изделием.
\item[Сеанс администратора] — интервал работы администратора с интерфейсом управления
  от успешной аутентификации до выхода или принудительного завершения.
\item[Предустановленный администратор] — учётная запись администратора,
  существующая в ПО МЭ с момента установки.
\item[Журнал событий] — хранилище записей о событиях, указанных в подразделе~\ref{sec:reg}.
\item[Журнал пакетов] — хранилище записей об обработанных пакетах.
\item[Подсистема] — функционально законченная часть ПО МЭ, допускающая
  независимый запуск, остановку и перезапуск.
\item[Контролируемые файлы] — исполняемые и конфигурационные файлы ПО МЭ,
  подлежащие контролю целостности.
\item[Режим аварийной поддержки] — режим функционирования ПО МЭ после выявления
  нарушения целостности, в котором обработка трафика заблокирована и доступны
  только функции восстановления.
\item[Служебный интерфейс] — виртуальный интерфейс ОС типа TAP, парный
  обслуживаемому интерфейсу, через который подсистема фильтрации доставляет
  кадры стеку ОС изделия и получает кадры, сформированные изделием (ОБР-05).
\item[Транзитный трафик] — кадры, принятые одним обслуживаемым интерфейсом
  и передаваемые на другой обслуживаемый интерфейс.
\item[ЧЗ, ТЧ] — уровни прав доступа администратора «чтение/запись» и «только чтение».
\item[CLI] — командный интерфейс управления (англ. \textenglish{Command Line Interface}).
\item[ОС] — операционная система.
\end{description}
\newpage

% ======================= НАЗНАЧЕНИЕ РАЗРАБОТКИ ========================
\section{Назначение разработки}

\subsection{Функциональное назначение}

Функциональным назначением ПО МЭ является:
\begin{itemize}
\item фильтрация сетевого трафика, поступающего на интерфейсы изделия,
  в соответствии с политикой доступа, в том числе с учётом состояния сеансов;
\item трансляция IPv4-адресов и TCP/UDP-портов источника и приёмника
      (Source NAT, Destination NAT, преобразование портов) по правилам,
      хранящимся отдельно от правил фильтрации;
\item регистрация событий безопасности и сведений об обработанных пакетах;
\item идентификация и аутентификация администраторов и разграничение их прав доступа;
\item локальное и удалённое управление изделием через командный и web-интерфейсы;
\item контроль целостности исполняемых и конфигурационных файлов ПО МЭ
  и блокирование трафика при её нарушении.
\end{itemize}

\subsection{Эксплуатационное назначение}

ПО МЭ эксплуатируется на изделии, устанавливаемом в разрыв между защищаемыми
сегментами сети (сегменты подключаются к разным интерфейсам изделия).
Эксплуатацию осуществляют администраторы безопасности, выполняющие
настройку политики доступа, контроль журналов и сопровождение изделия.
\newpage
% =================== ТРЕБОВАНИЯ К ПРОГРАММЕ ===========================
\section{Требования к программе}

\subsection{Требования к функциональным характеристикам}\label{sec:func}

Требования снабжены идентификаторами вида \textbf{ГРП-NN}, где ГРП~— код группы
требований, NN~— порядковый номер. После текста требования указан пункт задания
на курсовую работу, из которого оно следует; требования, уточняющие задание
и необходимые для однозначной реализации и проверки, помечены как производные.
Требования, содержащие слово «должно», являются обязательными.

% --------------------------- 4.1.1 Общие -----------------------------
\subsubsection{Общие требования (ОБЩ)}

\begin{reqs}
\item[ОБЩ-01] ПО МЭ должно функционировать под управлением ОС семейства Linux
  (ядро версии 5.10 и выше) на аппаратной платформе с архитектурой x86-64.
  \src{1.1}
\item[ОБЩ-02] ПО МЭ должно поддерживать от 1 до 6 сетевых интерфейсов Ethernet
  со скоростью 1~Гбит/с (1000BASE-T). Количество и перечень обслуживаемых
  интерфейсов определяется конфигурацией ПО МЭ. \src{1.2}
\item[ОБЩ-03] Обслуживаемые интерфейсы должны обозначаться в правилах, журналах
  и интерфейсах управления едиными логическими именами (\texttt{port1}~\ldots~\texttt{port6}),
  однозначно сопоставленными системным именам сетевых интерфейсов ОС. \der
\item[ОБЩ-04] В составе ПО МЭ должны быть выделены как минимум следующие подсистемы,
  допускающие независимый запуск, остановку и перезапуск:
  \begin{itemize}
  \item подсистема фильтрации (включая трансляцию адресов и контроль состояний);
  \item подсистема регистрации;
  \item подсистема управления (командный и web-интерфейсы);
  \item подсистема контроля целостности.
  \end{itemize}
  Окончательный состав подсистем определяется на этапе разработки архитектуры.
  \src{3.1, 5.11}
\end{reqs}

% ------------------------ 4.1.2 Фильтрация ---------------------------
\subsubsection{Требования к контролю и фильтрации трафика (ФЛТ)}

\begin{reqs}
\item[ФЛТ-01] ПО МЭ должно обрабатывать весь трафик, поступающий на обслуживаемые
  интерфейсы, в соответствии с политикой доступа. Передача пакета с одного
  интерфейса на другой в обход политики доступа не допускается. \src{2.1}
\item[ФЛТ-02] Политика доступа должна представлять собой список правил,
  пронумерованных уникальными целыми положительными числами. Правила просматриваются
  в порядке возрастания номеров; к пакету применяется первое активное правило,
  всем критериям которого пакет соответствует. \src{2.1, 2.3}
\item[ФЛТ-03] Пакет, к которому не применилось ни одно правило, должен удаляться. Регистрация таких пакетов
  в журнале пакетов должна включаться и отключаться параметром конфигурации
  (по умолчанию~— разрешена); в записи журнала указывается номер правила~0. \der
\item[ФЛТ-04] Должна обеспечиваться фильтрация по следующим критериям, каждый
  из которых в правиле может иметь значение «любой»: \src{2.2}
  \begin{enumerate}
  \item входной интерфейс~— один или несколько логических имён интерфейсов;
  \item MAC-адрес источника и MAC-адрес приёмника;
  \item IPv4-адрес источника и IPv4-адрес приёмника~— отдельный адрес или
    подсеть в нотации CIDR (например, \texttt{192.168.1.0/24});
  \item протокол, инкапсулированный в IP-датаграмму,~— номер протокола (0--255)
    или мнемоническое обозначение (\texttt{tcp}, \texttt{udp}, \texttt{icmp});
  \item TCP/UDP-порт источника и TCP/UDP-порт приёмника~— отдельный номер порта
    или диапазон (1--65535);
  \item тип и код ICMPv4 (0--255 каждый).
  \end{enumerate}
\item[ФЛТ-05] Критерии по портам допускаются только в правилах с протоколом
  TCP или UDP, критерии по типу и коду ICMPv4~— только в правилах с протоколом ICMP.
  Правило, нарушающее данное условие или содержащее синтаксически неверные
  значения, не должно приниматься с выдачей администратору сообщения об ошибке. \der
\item[ФЛТ-06] Кадры Ethernet, не содержащие IPv4-датаграмм (например, ARP),
  должны обрабатываться политикой доступа; для таких кадров правило может быть
  применено, только если все критерии сетевого и транспортного уровней в нём
  имеют значение «любой». \der
\item[ФЛТ-07] Правило фильтрации должно содержать следующие атрибуты: \src{2.3}
  \begin{enumerate}
  \item номер правила;
  \item критерии отбора пакетов (ФЛТ-04);
  \item действие: «передать» с указанием одного или нескольких выходных
    интерфейсов; «перейти» с указанием номера правила; «удалить»;
  \item контроль состояния: «контролировать» / «не контролировать»;
  \item активность: «активно» / «неактивно»; неактивное правило при обработке
    пакетов не применяется;
  \item флаг регистрации: «регистрировать» / «не регистрировать» пакет в журнале пакетов;
  \item комментарий~— произвольная строка длиной до 255 символов;
  \end{enumerate}
\item[ФЛТ-08] При выполнении действия «передать» пакет должен передаваться на
  указанные в правиле выходные интерфейсы. Если выходной интерфейс выключен,
  пакет удаляется. \src{2.3}

  \textit{Примечание~— Способ определения канального адреса следующего узла
  (прозрачная передача на канальном уровне или маршрутизация с использованием
  системных маршрутов) определяется на этапе разработки архитектуры.}
\item[ФЛТ-09] При выполнении действия «перейти» просмотр политики должен
  продолжаться с правила с наименьшим номером, большим или равным указанному.
  Номер правила перехода должен быть больше номера текущего правила (исключение
  зацикливания); правило, нарушающее это условие, не принимается. Если правил
  с номером, большим или равным указанному, нет, применяется ФЛТ-03. \src{2.3}

\item[ФЛТ-10] Трансляция адресов и портов не является атрибутом правила
фильтрации и выполняется по отдельным правилам NAT в соответствии с
требованиями НАТ-01 -- НАТ-15. Правило фильтрации не должно содержать
параметров трансляции. \src{2.4}


\item[ФЛТ-11]Обратная трансляция ответных пакетов сеанса выполняется по
записи таблицы состояний в соответствии с требованием НАТ-12.
(производное требование)
\end{reqs}

% ----------------------- 4.1.3 Контроль состояний --------------------
\subsubsection{Требования к контролю состояний (СОС)}

\begin{reqs}
\item[СОС-01] Для пакета, пропущенного правилом с атрибутом «контролировать»,
  ПО МЭ должно создавать запись в таблице состояний, если пакет инициирует новый
  сеанс. Для TCP новый сеанс может быть инициирован только сегментом с флагом SYN
  без флага ACK; прочие TCP-сегменты, не принадлежащие существующему сеансу,
  при обработке таким правилом удаляются. \src{2.3}
\item[СОС-02]Пакеты, принадлежащие сеансу, для которого существует запись
в таблице состояний (в прямом и обратном направлениях), должны
пропускаться без повторного просмотра политики доступа и таблиц NAT, с
выполнением прямой и обратной трансляции по значениям, сохранённым в
записи. Такие пакеты регистрируются в журнале пакетов с номером правила
фильтрации, создавшего запись, если у этого правила установлен флаг
регистрации. (п. 2.3 задания)
\item[СОС-03] Запись таблицы состояний должна содержать: протокол;
исходные (до трансляции) и транслированные IP-адреса; исходные и
транслированные порты (для TCP/UDP) или идентификатор (для ICMP); входной
и выходной интерфейсы; состояние сеанса (для TCP); номер создавшего
правила фильтрации; номера применённых правил DNAT и SNAT (0, если
трансляция не применялась); время создания; оставшееся время жизни.
(п. 5.7 задания)
\item[СОС-04] Записи должны удаляться по истечении тайм-аута неактивности
  и при штатном завершении TCP-соединения. Должны поддерживаться настраиваемые
  параметры (значения по умолчанию указаны в скобках): тайм-аут установленного
  TCP-сеанса (3600~с), тайм-аут TCP-сеанса в прочих состояниях (120~с), отсутствие ответа на SYN-пакет (5~с),
  тайм-аут UDP (60~с), тайм-аут ICMP (30~с), максимальное число записей (65536).
  При заполнении таблицы новые сеансы не создаются, а инициирующие их пакеты
  удаляются. \der
\item[СОС-05] Механизм контроля состояний должен включаться и отключаться
  администратором. При отключении атрибут «контролировать» в правилах
  игнорируется, а таблица состояний очищается. \src{5.7}
\end{reqs}

% -------------------------- 4.1.4 Регистрация ------------------------
\subsubsection{Требования к регистрации (РЕГ)}\label{sec:reg}

\begin{reqs}
\item[РЕГ-01] Должна обеспечиваться регистрация и хранение информации о следующих
  событиях: \src{3.1}
  \begin{itemize}
  \item вход и выход администратора (включая неуспешные попытки входа
    и завершение сеанса по тайм-ауту или принудительно);
  \item включение и выключение МЭ и его подсистем (включая аварийное
    завершение подсистемы);
  \item действия администратора по изменению политики доступа;
  \item действия администратора по изменению конфигурационных параметров
    в процессе управления МЭ;
  \item события контроля целостности (выявление нарушения, переход в режим
    аварийной поддержки и выход из него).
  \end{itemize}
  Перечень кодов событий приведён в приложении~А.
\item[РЕГ-02] Запись журнала событий должна содержать: \src{3.2}
  \begin{itemize}
  \item дату и время события;
  \item код и описание события;
  \item идентификатор администратора, действия которого привели к регистрации
    события, либо наименование подсистемы~— инициатора события;
  \item IP-адрес компьютера администратора в случае удалённого администрирования.
  \item успешность или неуспешность события
  \end{itemize}
\item[РЕГ-03] Для пакетов, обработанных правилом с установленным флагом регистрации,
  в журнале пакетов должна сохраняться следующая информация: \src{3.3}
  \begin{itemize}
  \item дата и время регистрации (с точностью до миллисекунды);
  \item действие, произведённое над пакетом («передан», «удалён»), и вид
      выполненной трансляции (нет, DNAT, SNAT, DNAT и SNAT);
\item номера применённых правил DNAT и SNAT (0, если трансляция не
      применялась);
\item IP-адреса и порты отправителя и получателя до и после трансляции;
  \item номер правила, которым обработан пакет;
  \item входной и выходной (выходные) интерфейсы;
  \item протоколы, передаваемые в пакете (стек протоколов, например
  \item состояние соединения (установлено/не установлено/прервано);
    \texttt{Ethernet/IPv4/TCP});
  \item MAC-адреса отправителя и получателя;
  \item IP-адреса отправителя и получателя;
  \item номера портов отправителя и получателя для TCP/UDP;
  \item тип и код для ICMPv4.
  \end{itemize}
\item[РЕГ-04] Журналы должны храниться в энергонезависимой памяти и сохраняться
  при перезагрузке и выключении изделия. \src{3.1, 3.3}
\item[РЕГ-05] Журнал событий не должен изменяться и очищаться средствами
  интерфейсов управления. Очистка журнала пакетов допускается только по команде
  администратора с правами ЧЗ и регистрируется в журнале событий. \src{5.6}
\item[РЕГ-06] Для каждого журнала должен задаваться максимальный размер. При его
  достижении удаляются наиболее старые записи, о чём регистрируется событие. \der
\end{reqs}

% ------------------ 4.1.5 Идентификация и аутентификация -------------
\subsubsection{Требования к идентификации, аутентификации и разграничению доступа (ДСТ)}

\begin{reqs}
\item[ДСТ-01] Доступ к функциям управления через любой интерфейс управления
  (локальный и удалённый CLI, web-интерфейс) должен предоставляться только после
  аутентификации администратора по идентификатору (имени учётной записи) и паролю.
  До успешной аутентификации функции управления должны быть недоступны. \src{4.1}
\item[ДСТ-02] Пароли должны храниться только в виде значений криптографической
  хеш-функции с использованием «соли». Пароль не должен отображаться при вводе
  и выводиться в каких-либо интерфейсах и журналах. Минимальная длина пароля~—
  8 символов. \der
\item[ДСТ-03] Должны поддерживаться два уровня прав доступа: \src{4.2}
  \begin{itemize}
  \item «чтение/запись» (ЧЗ)~— доступ ко всем функциям управления;
  \item «только чтение» (ТЧ)~— доступ только к функциям просмотра.
  \end{itemize}
  Распределение функций по уровням прав приведено в приложении~Б. Попытка
  выполнения функции, на которую у администратора нет прав, должна отклоняться
  с выдачей сообщения и регистрироваться в журнале событий.
\item[ДСТ-04] В ПО МЭ должна присутствовать предустановленная учётная запись
  администратора (идентификатор \texttt{admin}) с правами ЧЗ и предустановленным
  паролем, указанным в эксплуатационной документации. Попытка удаления данной
  учётной записи или понижения её прав должна блокироваться и регистрироваться
  в журнале событий. Изменение пароля предустановленного администратора
  допускается. \src{4.3}
\item[ДСТ-05] Если на момент аутентификации администратора с правами ЧЗ в системе
  уже существует активный сеанс с правами ЧЗ (в том числе другого сеанса той же
  учётной записи), права нового сеанса должны понижаться до ТЧ. Администратору
  выдаётся сообщение о понижении прав, событие регистрируется в журнале событий.
  Повышение прав сеанса после завершения более раннего сеанса не производится~—
  для получения прав ЧЗ требуется повторный вход. \src{4.4}
\item[ДСТ-06] Сеанс администратора должен завершаться автоматически после периода
  неактивности, задаваемого параметром конфигурации (по умолчанию 15~мин). \der
\item[ДСТ-07] Удалённое управление должно осуществляться по защищённым каналам:
  CLI~— по протоколу SSH, web-интерфейс~— по протоколу HTTPS (TLS версии 1.2 и выше). \der
\end{reqs}

% --------------------------- 4.1.6 Управление ------------------------
\subsubsection{Требования к управлению (УПР)}

\begin{reqs}
\item[УПР-01] Должно обеспечиваться локальное управление (с консоли изделия)
  и удалённое управление (по сети через интерфейсы изделия, на которых
  настроен IP-адрес). \src{5.1}
\item[УПР-02] Должны быть реализованы командный интерфейс (CLI) и web-интерфейс.
  Каждый из них должен предоставлять полный перечень функций управления,
  приведённый в приложении~Б. \src{5.2}
\item[УПР-03] Должны обеспечиваться просмотр и редактирование правил
фильтрации, правил DNAT и правил SNAT: добавление, изменение, удаление,
перенумерация правил, изменение атрибутов (в том числе активности).
Операции выполняются через API управления (АПИ-01). \src{5.3}
\item[УПР-04] Должны обеспечиваться просмотр журнала событий и выборка записей
  по коду события и по интервалу времени регистрации. \src{5.4}
\item[УПР-05] Должны обеспечиваться просмотр журнала пакетов и выборка записей
  по следующим критериям: действие над пакетом; входной и выходной интерфейсы;
  номер правила; MAC-адреса отправителя и получателя; IP-адреса отправителя
  и получателя; порты отправителя и получателя; интервал времени регистрации. \src{5.5}
\item[УПР-06] При выборке записей журналов критерии должны комбинироваться
  по логическому «И»; незаданный критерий не ограничивает выборку. \der
\item[УПР-07] Должна обеспечиваться очистка журнала пакетов. \src{5.6}
\item[УПР-08] Должно обеспечиваться управление контролем состояний: просмотр
  настроек механизма; включение и отключение механизма; просмотр таблицы
  состояний; удаление отдельных записей и всех записей таблицы состояний. \src{5.7}
\item[УПР-09] Должно обеспечиваться управление конфигурациями: \src{5.8}
  \begin{itemize}
  \item сохранение текущей конфигурации;
  \item вывод текущей и сохранённой конфигурации;
  \item применение сохранённой конфигурации;
  \item применение конфигурации по умолчанию.
  \end{itemize}
\item[УПР-10] Конфигурация должна включать: политику доступа, таблицы DNAT и SNAT, параметры
  интерфейсов, системные маршруты, параметры контроля состояний, параметры
  регистрации, параметры контроля целостности, параметры сеансов администраторов.
  Учётные записи администраторов в конфигурацию не входят и при применении
  конфигурации не изменяются. Сохранённая конфигурация должна применяться
  автоматически при запуске ПО МЭ; при её отсутствии применяется конфигурация
  по умолчанию. Конфигурация должна храниться в текстовом формате. \der
\item[УПР-11] Конфигурация по умолчанию должна содержать пустую политику доступа
  (весь трафик удаляется по ФЛТ-03), пустые таблицы DNAT и SNAT, включённый механизм контроля состояний со
  значениями параметров по умолчанию и интерфейсы без назначенных IP-адресов. \der
\item[УПР-12] Должно обеспечиваться управление учётными записями
  администраторов: добавление, удаление, редактирование (смена пароля, уровня
  прав) учётной записи; просмотр базы учётных записей (без паролей); просмотр
  активных сеансов (идентификатор, интерфейс управления, IP-адрес, время входа,
  уровень прав); принудительный сброс активных сеансов. \src{5.9}
\item[УПР-13] Должна обеспечиваться настройка параметров Ethernet-интерфейсов:
  просмотр состояния (логическое и системное имя, MAC-адрес, состояние канала,
  скорость, административное состояние, IP-адрес и маска, счётчики пакетов);
  включение и выключение интерфейса; назначение и удаление IP-адреса
  интерфейса. \src{5.10}
\item[УПР-14] Должно обеспечиваться управление системными настройками
  и режимами функционирования: \src{5.11}
  \begin{itemize}
  \item просмотр и настройка даты и времени;
  \item добавление, удаление и просмотр системных маршрутов;
  \item выключение изделия;
  \item проверка целостности контролируемых файлов по запросу;
  \item перезагрузка изделия;
  \item перезагрузка любой из подсистем ПО МЭ;
  \item просмотр системной информации: процессор и его загрузка, объём
    и использование памяти, интерфейсы, состояние подсистем.
  \end{itemize}
\item[УПР-15] Операции выключения и перезагрузки изделия, применения
  конфигурации, очистки журнала пакетов и удаления учётной записи должны
  выполняться только после подтверждения администратором. \der
\end{reqs}

% ----------------------- 4.1.7 Контроль целостности ------------------
\subsubsection{Требования к контролю целостности (ЦЛС)}

\begin{reqs}
\item[ЦЛС-01] Контролю целостности должны подлежать все исполняемые файлы
  ПО МЭ и его конфигурационные файлы (включая сохранённые конфигурации
  и базу учётных записей). Контроль выполняется сравнением значений хеш-функции
  SHA-256 файлов с эталонными значениями. \src{6.1}
\item[ЦЛС-02] Эталонные значения должны формироваться при установке ПО МЭ
  и автоматически обновляться при легитимном изменении конфигурационных файлов
  средствами интерфейсов управления. \der
\item[ЦЛС-03] Контроль целостности должен выполняться: при запуске ПО МЭ;
  периодически в автоматическом режиме с периодом, задаваемым администратором
  (от 1 до 1440~мин, по умолчанию 10~мин); по запросу администратора. \src{6.1}
\item[ЦЛС-04] При выявлении нарушения целостности хотя бы одного контролируемого
  файла ПО МЭ должно: \src{6.2}
  \begin{itemize}
  \item заблокировать весь обрабатываемый трафик на всех интерфейсах;
  \item зарегистрировать событие с указанием файлов, целостность которых нарушена;
  \item перейти в режим аварийной поддержки.
  \end{itemize}
\item[ЦЛС-05] В режиме аварийной поддержки должно быть доступно локальное
  управление через CLI после аутентификации администратора с правами ЧЗ
  со следующими функциями: просмотр перечня файлов с нарушенной целостностью;
  восстановление файлов из резервной копии; повторная проверка целостности;
  просмотр журнала событий; выключение и перезагрузка изделия. \src{6.2}
\item[ЦЛС-06] Обработка трафика должна возобновляться автоматически после того,
  как повторная проверка подтвердит целостность всех контролируемых файлов.
  Выход из режима аварийной поддержки регистрируется в журнале событий. \src{6.2}
\item[ЦЛС-07] Если нарушение целостности выявлено при запуске, ПО МЭ должно
  сразу запускаться в режиме аварийной поддержки без пропуска трафика. \der
\end{reqs}


\subsubsection{Требования к трансляции адресов и портов (НАТ)}
\begin{reqs}
\item[НАТ-01] Правила фильтрации и трансляции адресов должны храниться в одном файле политики. Файл должен содержать упорядоченные правила фильтрации, DNAT и SNAT с указанием типа каждого правила. Правила NAT определяют преобразование пакета; решение о его передаче или удалении принимается правилами фильтрации. \src{2.4}
\item[НАТ-02] Каждый пакет обрабатывается независимо. Сначала к пакету применяется подходящее правило DNAT, затем политика фильтрации с учётом изменённого адреса приёмника. Если пакет разрешён, к нему применяется подходящее правило SNAT перед передачей. \der
\item[НАТ-03] Правила каждого типа проверяются в порядке их следования в файле. Применяется первое активное правило соответствующего типа, критериям которого удовлетворяет пакет. Если подходящего правила данного типа нет, соответствующий адрес пакета не изменяется. \der
\item[НАТ-04] Правило NAT должно содержать номер, тип (DNAT или SNAT), признак активности, критерии отбора и IPv4-адрес, на который заменяется соответствующий адрес пакета. Правило может содержать комментарий. \der
\item[НАТ-05] Критерии отбора правила NAT должны включать входной интерфейс, IPv4-адреса источника и приёмника, протокол IP и порты источника и приёмника для TCP/UDP. Неуказанный критерий не ограничивает отбор; указанные критерии проверяются совместно. \der
\item[НАТ-06] При совпадении пакета с правилом DNAT IPv4-адрес приёмника заменяется адресом, заданным в правиле. При совпадении с правилом SNAT IPv4-адрес источника заменяется адресом, заданным в правиле. \src{2.4}
\item[НАТ-07] Правило NAT может задавать замену порта источника для SNAT или порта приёмника для DNAT. Замена порта применяется только к пакетам TCP или UDP при наличии соответствующего транспортного заголовка. \der
\item[НАТ-08] При отсутствии параметра замены порта исходный номер порта сохраняется. Замена адреса выполняется независимо от наличия параметра порта. \der
\item[НАТ-9] При изменении заголовков должны обеспечиваться корректные контрольные суммы IPv4 и соответствующего транспортного протокола, если его контрольная сумма зависит от изменённых полей. \der
\item[НАТ-10] Изменённая политика должна применяться без перезапуска подсистемы фильтрации. При этом каждый пакет обрабатывается по целостной версии политики. \der
\item[НАТ-11] Правила NAT не должны самостоятельно разрешать передачу пакета, запрещённого политикой фильтрации. \der
\item[НАТ-12] При добавлении или изменении правила должны проверяться тип правила, корректность критериев, IPv4-адреса трансляции и совместимость заданного порта с протоколом. Некорректное правило не должно применяться. \der
\item[НАТ-13] Для обратного направления трафика трансляция выполняется только при наличии подходящих явно заданных правил. Автоматическое восстановление адресов по состоянию соединения не выполняется. \der
\end{reqs}

\subsubsection{Требования к структурам данных политики (СТР)}
\begin{reqs}
\item[СТР-01] В ПО МЭ политика доступа и правила NAT должны быть представлены
  тремя независимыми структурами данных: таблицей правил фильтрации, таблицей
  правил DNAT и таблицей правил SNAT. Каждая таблица имеет собственную
  нумерацию правил и собственные операции добавления, изменения, удаления и
  перенумерации. Изменение одной таблицы не должно изменять содержимое
  другой. Правило фильтрации не должно содержать полей трансляции, правило
  NAT --- полей действия фильтрации и регистрации. \der
\item[СТР-02] Правило фильтрации должно иметь состав полей, приведённый в
  таблице~\ref{tab:str-filter}. Значение «любой» представляется отсутствием
  поля (или значением \texttt{null}). \der
\end{reqs}
 

{\small
\begin{longtable}{|L{3.6cm}|L{3.3cm}|L{7.2cm}|}
\caption{Структура правила фильтрации}\label{tab:str-filter}\\
\hline
\textbf{Поле} & \textbf{Тип} & \textbf{Допустимые значения и ограничения}\\
\hline\endfirsthead
\hline
\textbf{Поле} & \textbf{Тип} & \textbf{Допустимые значения и ограничения}\\
\hline\endhead
\texttt{id} & целое & $>0$, уникально в таблице фильтрации\\ \hline
\texttt{active} & логическое & активно / неактивно\\ \hline
\texttt{in\_ifaces} & список имён & \texttt{port1} \dots \texttt{port6} либо «любой»\\ \hline
\texttt{src\_mac}, \texttt{dst\_mac} & MAC-адрес & либо «любой»\\ \hline
\texttt{src\_ip}, \texttt{dst\_ip} & IPv4 / CIDR & либо «любой»\\ \hline
\texttt{proto} & число или имя & 0--255, \texttt{tcp}, \texttt{udp}, \texttt{icmp} либо «любой»\\ \hline
\texttt{src\_port}, \texttt{dst\_port} & порт или диапазон & 1--65535; только при \texttt{tcp}/\texttt{udp}\\ \hline
\texttt{icmp\_type}, \texttt{icmp\_code} & целое & 0--255; только при \texttt{icmp}\\ \hline
\texttt{action} & перечисление & \texttt{forward}, \texttt{jump}, \texttt{drop}\\ \hline
\texttt{out\_ifaces} & список имён & обязателен и не пуст при \texttt{forward}\\ \hline
\texttt{jump\_to} & целое & обязателен при \texttt{jump}, строго больше \texttt{id} (ФЛТ-09)\\ \hline
\texttt{stateful} & логическое & контролировать / не контролировать\\ \hline
\texttt{log} & логическое & регистрировать / не регистрировать\\ \hline
\texttt{comment} & строка & до 255 символов\\ \hline
\end{longtable}
}
 
\begin{reqs}
\item[СТР-03] Правило DNAT должно иметь состав полей, приведённый в
  таблице~\ref{tab:str-dnat}. \der
\end{reqs}
 
{\small
\begin{longtable}{|L{3.6cm}|L{3.3cm}|L{7.2cm}|}
\caption{Структура правила DNAT}\label{tab:str-dnat}\\
\hline
\textbf{Поле} & \textbf{Тип} & \textbf{Допустимые значения и ограничения}\\
\hline\endfirsthead
\hline
\textbf{Поле} & \textbf{Тип} & \textbf{Допустимые значения и ограничения}\\
\hline\endhead
\texttt{id} & целое & $>0$, уникально в таблице DNAT\\ \hline
\texttt{active} & логическое & активно / неактивно\\ \hline
\texttt{in\_ifaces} & список имён & \texttt{port1} \dots \texttt{port6} либо «любой»\\ \hline
\texttt{src\_ip}, \texttt{dst\_ip} & IPv4 / CIDR & либо «любой»\\ \hline
\texttt{proto} & число или имя & как в таблице~\ref{tab:str-filter}\\ \hline
\texttt{src\_port}, \texttt{dst\_port} & порт или диапазон & только при \texttt{tcp}/\texttt{udp}\\ \hline
\texttt{action} & перечисление & \texttt{translate}, \texttt{skip}\\ \hline
\texttt{to\_ip} & IPv4 & новый адрес приёмника; необязателен\\ \hline
\texttt{to\_port} & порт или диапазон & новый порт приёмника; только при \texttt{tcp}/\texttt{udp}; необязателен\\ \hline
\texttt{comment} & строка & до 255 символов\\ \hline
\end{longtable}
}
 
\begin{reqs}
\item[СТР-04] Правило SNAT должно иметь состав полей, приведённый в
  таблице~\ref{tab:str-snat}. \der
\end{reqs}
{\small
\begin{longtable}{|L{3.6cm}|L{3.3cm}|L{7.2cm}|}
\caption{Структура правила SNAT}\label{tab:str-snat}\\
\hline
\textbf{Поле} & \textbf{Тип} & \textbf{Допустимые значения и ограничения}\\
\hline\endfirsthead
\hline
\textbf{Поле} & \textbf{Тип} & \textbf{Допустимые значения и ограничения}\\
\hline\endhead
\texttt{id} & целое & $>0$, уникально в таблице SNAT\\ \hline
\texttt{active} & логическое & активно / неактивно\\ \hline
\texttt{in\_ifaces}, \texttt{out\_ifaces} & списки имён & \texttt{port1} \dots \texttt{port6} либо «любой»\\ \hline
\texttt{src\_ip}, \texttt{dst\_ip} & IPv4 / CIDR & либо «любой»\\ \hline
\texttt{proto} & число или имя & как в таблице~\ref{tab:str-filter}\\ \hline
\texttt{src\_port}, \texttt{dst\_port} & порт или диапазон & только при \texttt{tcp}/\texttt{udp}\\ \hline
\texttt{action} & перечисление & \texttt{translate}, \texttt{skip}\\ \hline
\texttt{to\_ip} & IPv4 или \texttt{iface} & новый адрес источника либо адрес выходного интерфейса; необязателен\\ \hline
\texttt{to\_port} & порт или диапазон & новый порт источника (диапазон --- для динамического выбора, НАТ-09); необязателен\\ \hline
\texttt{comment} & строка & до 255 символов\\ \hline
\end{longtable}
}
 
\begin{reqs}
\item[СТР-05] Критерии отбора пакетов (интерфейсы, адреса, протокол, порты) во
  всех трёх таблицах должны использовать одни и те же типы значений и одну и
  ту же процедуру проверки. Это исключает расхождение форматов между
  фильтрацией и NAT. \der
\item[СТР-06] В текстовой конфигурации (УПР-10) каждая таблица должна храниться
  в отдельной секции (\texttt{filter}, \texttt{dnat}, \texttt{snat}).
  Применение конфигурации заменяет все три таблицы. \der
\item[СТР-07] Замена таблицы при изменении должна быть атомарной: каждый пакет
  обрабатывается либо по прежней, либо по новой версии таблицы целиком, без
  промежуточных состояний. Это обеспечивает выполнение ФЛТ-12 и НАТ-13.
  \der
\item[СТР-08] Записи таблицы состояний должны ссылаться на правила только по
  номеру и таблице (фильтрация, DNAT, SNAT), а не хранить копию правила.
  \der
\end{reqs}
 
 
\subsubsection{Требования к API управления политикой и NAT (АПИ)}
\begin{reqs}
\item[АПИ-01] Подсистема фильтрации должна предоставлять локальный программный интерфейс управления единой политикой, включающей правила фильтрации, DNAT и SNAT. Интерфейс предназначен для обмена между подсистемой фильтрации и программами управления. \der
\item[АПИ-02] Обмен с подсистемой фильтрации должен выполняться через локальный потоковый сокет семейства UNIX. Веб-сервер выдаёт страницы браузеру, принимает данные пользователя и передаёт соответствующие команды в подсистему фильтрации через этот сокет. HTTPS между браузером и веб-сервером и SSH для удалённого командного доступа не являются протоколом этого локального API. \der
\item[АПИ-03] Запросы и ответы должны передаваться в кодировке UTF-8 в формате JSON: один объект на строку, завершённую символом перевода строки. Строки внутри значений JSON должны экранироваться по правилам JSON. На каждый корректно принятый запрос формируется один ответ. \der
\item[АПИ-04] Запрос должен содержать поля \texttt{id} (идентификатор запроса), \texttt{command} (имя команды) и \texttt{params} (объект параметров, пустой при их отсутствии). Ответ должен содержать тот же \texttt{id} и поле \texttt{status} со значением \texttt{ok} или \texttt{error}; при успешном выполнении результат передаётся в поле \texttt{result}, при ошибке --- код и описание в объекте \texttt{error}. Если идентификатор не удалось извлечь, ответ содержит \texttt{id: null}. \der
\item[АПИ-05] API должен поддерживать команды: \texttt{list\_rules} --- просмотр упорядоченных правил с необязательным отбором по типу; \texttt{get\_rule} --- просмотр правила по идентификатору; \texttt{add\_rule}, \texttt{update\_rule}, \texttt{delete\_rule} --- добавление, изменение и удаление; \texttt{move\_rule} --- изменение положения; \texttt{set\_enabled} --- изменение активности; \texttt{get\_policy} --- получение всей политики. Команды применяются к правилам фильтрации и NAT соответствующего типа. \der
\item[АПИ-06] Параметры операций над отдельным правилом должны содержать его идентификатор; для добавления и изменения передаётся объект правила с типом, критериями и параметрами действия. Для правила NAT передаются тип SNAT или DNAT и новый IPv4-адрес, при необходимости --- новый порт. Для перемещения указывается позиция в порядке правил соответствующего типа, для изменения активности --- логическое значение. При добавлении идентификатор назначается подсистемой и возвращается в ответе. \der
\item[АПИ-07] Ответ на просмотр должен содержать запрошенные правила со всеми их полями в порядке применения и номер текущей версии политики. Ответ на успешное изменение должен содержать идентификатор изменённого правила и новую версию политики. \der
\item[АПИ-08] Подсистема должна проверять формат запроса, наличие и допустимость полей и параметров правила. Для каждого изменения клиент передаёт ожидаемую версию политики; при её несовпадении изменение отклоняется. \der
\item[АПИ-09] Ответ об ошибке должен содержать машинный код и понятное описание причины. Должны различаться как минимум ошибки некорректного запроса, неизвестной команды, неверных параметров, отсутствующего правила, конфликта версии, недостаточных прав и внутренней ошибки. При ошибке политика не изменяется. \der
\item[АПИ-10] Каждое изменение должно выполняться атомарно: после успешной проверки новая версия сохраняется в единый файл политики и применяется подсистемой фильтрации целиком; при неудаче сохраняется прежняя действующая версия. Успешный ответ возвращается после сохранения и применения. \der
\item[АПИ-11] Доступ к локальному сокету должен ограничиваться средствами ОС уполномоченными процессами. Веб-сервер и командный интерфейс должны проверять полномочия администратора до передачи команды: роль ТЧ допускает только просмотр, роль ЧЗ --- также изменение. Идентификатор администратора и источник обращения передаются для регистрации изменений. \der
\item[АПИ-12] Успешные операции изменения политики должны регистрироваться в журнале событий с идентификатором администратора, видом операции, идентификатором правила и временем. Интерфейс должен возвращать результат выполнения вызывающей стороне. \der
\end{reqs}

\subsubsection{Требования к реализации обработки Ethernet-кадров и IP-пакетов в ОС Linux (ОБР)}\label{sec:impl}

Требования данной группы определяют, каким образом ПО МЭ получает, обрабатывает
и передаёт кадры в ОС Linux, чтобы выполнение ФЛТ-01 и НДЖ-02 не зависело
от настроек сетевого стека ядра. Решение о пропуске каждого кадра принимает
только подсистема фильтрации, работающая в пространстве пользователя. Ядро ОС
используется как драйвер сетевых адаптеров и как сетевой стек собственных
служб изделия (SSH, HTTPS, API управления). Требования группы являются
производными и закрывают вопрос, оставленный открытым в примечании к ФЛТ-08.

\begin{reqs}
\item[ОБР-01] Подсистема фильтрации должна принимать и передавать кадры
  обслуживаемых интерфейсов через сокеты \texttt{AF\_PACKET} типа
  \texttt{SOCK\_RAW} с протоколом \texttt{ETH\_P\_ALL}: по одному сокету на
  интерфейс, привязанному к нему вызовом \texttt{bind()}. Приём выполняется
  через кольцевой буфер в разделяемой памяти (\texttt{PACKET\_RX\_RING},
  версия \texttt{TPACKET\_V3}), передача~— вызовом \texttt{send()} или через
  \texttt{PACKET\_TX\_RING}. Пересылка транзитного трафика средствами ядра
  (мост, маршрутизация ядра, решения netfilter/nftables по политике) не
  допускается. \der

  \textit{Примечание~— Допускается замена \texttt{AF\_PACKET} на
  \texttt{AF\_XDP}, если при этом выполняются все требования группы ОБР.}

\item[ОБР-02] Для исключения обхода политики доступа системный модуль должен
  до запуска подсистемы фильтрации привести ОС в состояние, указанное
  в таблице~\ref{tab:os-settings}. Правило \texttt{drop} в цепочке
  \texttt{ingress} семейства \texttt{netdev} не допускает обработку кадров
  обслуживаемых интерфейсов стеком ядра. Сокеты \texttt{AF\_PACKET} получают
  копию кадра раньше этой цепочки, поэтому подсистема фильтрации видит весь
  трафик. Правило сохраняется при остановке и аварийном завершении подсистемы
  фильтрации, что обеспечивает отказ в «закрытое» состояние (НДЖ-02).
  Соответствие состояния ОС таблице~\ref{tab:os-settings} проверяется при
  каждой проверке целостности (ЦЛС-03); расхождение рассматривается как
  нарушение целостности (ЦЛС-04). \der
\end{reqs}

{\small
\begin{longtable}{|L{5.4cm}|L{3.4cm}|L{5.3cm}|}
\caption{Настройки ОС, устанавливаемые системным модулем}\label{tab:os-settings}\\
\hline
\textbf{Параметр} & \textbf{Значение} & \textbf{Назначение}\\
\hline\endfirsthead
\hline
\textbf{Параметр} & \textbf{Значение} & \textbf{Назначение}\\
\hline\endhead
\texttt{net.ipv4.ip\_forward} & 0 & запрет пересылки IPv4 ядром\\ \hline
\texttt{net.\allowbreak ipv6.\allowbreak conf.\allowbreak all.\allowbreak forwarding} & 0 & запрет пересылки IPv6 ядром\\ \hline
\texttt{net.\allowbreak ipv6.\allowbreak conf.\allowbreak <if>.\allowbreak disable\_ipv6} & 1 для обслуживаемых и служебных интерфейсов & IPv6 не обрабатывается (ФЛТ-04 предусматривает только IPv4)\\ \hline
IP-адреса обслуживаемых интерфейсов & не назначаются & адреса изделия задаются только на служебных интерфейсах (ОБР-05)\\ \hline
Мосты, bond, VLAN-интерфейсы & обслуживаемые интерфейсы не включаются & исключение пересылки ядром\\ \hline
Таблица \texttt{netdev}, цепочка \texttt{ingress} & \texttt{drop} для всех кадров каждого обслуживаемого интерфейса & отказ в «закрытое» состояние\\ \hline
Режим интерфейса & \texttt{promisc on}, MTU 1500 & приём кадров с любым MAC-адресом приёмника\\ \hline
Аппаратные функции (\texttt{ethtool -K}) & \texttt{gro}, \texttt{lro}, \texttt{gso}, \texttt{tso}, \texttt{rxvlan}, \texttt{txvlan}~— \texttt{off} & кадры поступают в том виде, в каком переданы по линии, не длиннее 1522 байт\\ \hline
\end{longtable}
}

\begin{reqs}
\item[ОБР-03] Соответствие логических имён \texttt{port1}\ldots\texttt{port6}
  системным именам интерфейсов (ОБЩ-03) должно задаваться в конфигурации.
  Интерфейс, отсутствующий в ОС, при запуске помечается как недоступный,
  о чём регистрируется событие; работа с остальными интерфейсами
  продолжается. \der

\item[ОБР-04] Кадры, полученные сокетом с признаком
  \texttt{sll\_pkttype}~=\allowbreak{} \texttt{PACKET\_OUTGOING} (собственные переданные кадры
  изделия), должны игнорироваться, чтобы исключить их повторную обработку. \der

\item[ОБР-05] Для доставки трафика собственным службам изделия (УПР-01) каждому
  обслуживаемому интерфейсу \texttt{portN} должен соответствовать служебный
  интерфейс типа TAP с тем же MAC-адресом. IP-адрес, назначаемый
  интерфейсу \texttt{portN} по УПР-13, фактически назначается его служебному
  интерфейсу; системные маршруты (УПР-14) и таблица ARP ядра также относятся
  к служебным интерфейсам. Для политики доступа служебные интерфейсы
  представлены одним логическим именем \texttt{local}, которое допускается:
  \begin{itemize}
  \item в списке выходных интерфейсов действия «передать»~— кадр записывается
    в служебный интерфейс, парный входному порту, и доставляется стеку ОС;
  \item в критерии «входной интерфейс»~— для кадров, прочитанных из служебного
    интерфейса (трафик, сформированный изделием). Такой кадр передаётся только
    на парный порт и только если этот порт указан в списке выходных интерфейсов
    применённого правила; в противном случае кадр удаляется.
  \end{itemize}
  Таким образом, трафик управления и ARP-обмен изделия также проходят политику
  доступа (ФЛТ-01). При конфигурации по умолчанию (УПР-11) удалённое
  управление недоступно до добавления разрешающих правил через локальную
  консоль. \der

\item[ОБР-06] Разбор кадра должен выполняться в следующем порядке:
  \begin{enumerate}
  \item заголовок Ethernet~II; кадры с полем длины IEEE~802.3 (значение меньше
    \texttt{0x0600}) и кадры с меткой IEEE~802.1Q (\texttt{0x8100})
    обрабатываются как кадры без IPv4-датаграммы (ФЛТ-06);
  \item при EtherType \texttt{0x0800}~— заголовок IPv4: версия~4,
    $\mathrm{IHL} \ge 5$, общая длина не меньше длины заголовка и не больше
    длины полезной нагрузки кадра, верная контрольная сумма заголовка;
    опции IP пропускаются без анализа; байты дополнения кадра Ethernet
    за пределами общей длины датаграммы отбрасываются;
  \item заголовок транспортного уровня: TCP~— смещение данных не меньше~5
    и заголовок помещается в датаграмму; UDP~— поле длины согласовано
    с длиной датаграммы; ICMPv4~— не менее 8 байт.
  \end{enumerate}
  Кадр, не прошедший проверки, удаляется до просмотра политики и учитывается
  в счётчике ошибок разбора интерфейса (УПР-13). В журнал пакетов такие кадры
  не записываются, чтобы исключить переполнение журнала некорректным
  трафиком. Контрольные суммы TCP и UDP не проверяются; при трансляции они
  корректируются по RFC~1624, нулевая контрольная сумма UDP остаётся
  нулевой (уточнение НАТ-11). \der

\item[ОБР-07] Фрагментированные IPv4-датаграммы должны собираться до
  просмотра таблиц NAT и политики доступа по ключу (адрес источника, адрес
  приёмника, протокол, идентификатор). Должны поддерживаться: тайм-аут сборки
  30~с; ограничение числа одновременно собираемых датаграмм (по умолчанию
  1024); удаление датаграммы целиком при перекрытии фрагментов или при
  превышении итоговой длины 65535 байт. Если после обработки длина датаграммы
  превышает MTU выходного интерфейса, она фрагментируется по RFC~791, а при
  установленном флаге DF удаляется. Собранная датаграмма регистрируется
  в журнале пакетов одной записью. \der

\item[ОБР-08] Канальные адреса передаваемого кадра должны определяться
  следующим образом:
  \begin{enumerate}
  \item если MAC-адрес приёмника кадра не совпадает с MAC-адресом входного
    порта (в том числе широковещательный и групповой), кадр передаётся
    прозрачно, MAC-адреса не изменяются;
  \item если MAC-адрес приёмника совпадает с MAC-адресом входного порта,
    а выходной интерфейс отличен от \texttt{local}, выполняется
    маршрутизация: по таблице маршрутов ядра (запрос \texttt{RTM\_GETROUTE}
    через netlink) для адреса приёмника после DNAT определяется следующий
    узел (шлюз маршрута, если маршрут ведёт через служебный интерфейс
    выходного порта; иначе сам адрес приёмника); MAC-адрес источника
    заменяется MAC-адресом выходного порта, MAC-адрес приёмника~— адресом
    следующего узла; TTL уменьшается на~1, датаграмма с TTL~$\le 1$
    удаляется;
  \item MAC-адрес следующего узла берётся из\linebreak[2] ARP-кэша подсистемы фильтрации;
    при его отсутствии подсистема передаёт ARP-запрос с выходного порта,
    а пакет помещается в очередь ожидания (не более 3 пакетов на адрес,
    не более 1~с) и удаляется при отсутствии ответа. ARP-ответы, адресованные
    изделию, обновляют кэш до просмотра политики, но к стеку ОС
    передаются только по правилу политики (ОБР-05). Время жизни записи
    кэша~— 60~с;
  \item при отсутствии маршрута пакет удаляется.
  \end{enumerate}
  \der

\item[ОБР-09] Обработка каждого принятого кадра должна выполняться в
  следующем порядке:
  \begin{enumerate}
  \item получение кадра и метки времени из кольцевого буфера
    (\texttt{TPACKET\_V3} даёт точность не хуже 1~мс, что достаточно
    для РЕГ-03);
  \item удаление кадра, если ПО МЭ находится в режиме аварийной поддержки
    (ЦЛС-04);
  \item разбор и проверка заголовков (ОБР-06), сборка фрагментов (ОБР-07);
  \item поиск сеанса в таблице состояний; при наличии~— трансляция по записи
    (СОС-02, НАТ-12) и переход к п.~7;
  \item просмотр таблицы DNAT (НАТ-02, п.~1);
  \item просмотр политики доступа (ФЛТ-02--ФЛТ-09); при действии «удалить»~—
    переход к п.~9;
  \item для каждого выходного интерфейса: проверка его состояния (ФЛТ-08),
    просмотр таблицы SNAT (новый сеанс), создание записи таблицы состояний
    (СОС-01, НАТ-10), определение канальных адресов (ОБР-08), корректировка
    контрольных сумм (ОБР-06), фрагментация (ОБР-07);
  \item передача копии кадра в сокет каждого выходного интерфейса;
  \item формирование записи журнала пакетов, если она требуется (РЕГ-03).
  \end{enumerate}
  \der

\item[ОБР-10] Регистрация пакетов не должна задерживать обработку трафика.
  Записи журнала пакетов передаются подсистеме регистрации через
  буфер в памяти; при его переполнении записи отбрасываются, а число
  потерянных записей отображается в системной информации (УПР-14). \der

\item[ОБР-11] Состояние обслуживаемых интерфейсов (административное
  состояние, наличие несущей) должно отслеживаться по уведомлениям netlink
  (группа \texttt{RTNLGRP\_LINK}). Подсистема фильтрации должна вести
  собственные счётчики для каждого интерфейса: принято, передано, удалено
  политикой, ошибки разбора, ошибки передачи. Эти счётчики выводятся по
  УПР-13, так как системные счётчики ОС не учитывают решений политики. \der

\item[ОБР-12] Процесс подсистемы фильтрации должен работать без прав
  суперпользователя, с возможностями \texttt{CAP\_NET\_RAW} и
  \texttt{CAP\_NET\_ADMIN}. Изменение таблиц правил и таблицы состояний
  должно быть безопасным при многопоточной обработке; замена таблиц
  выполняется атомарно (СТР-07). \der

\item[ОБР-13] При политике доступа до 100 правил и таблицах DNAT и SNAT
  до 100 правил каждая ПО МЭ должно обеспечивать передачу трафика между
  двумя физическими интерфейсами 1~Гбит/с со скоростью не менее 100~Мбит/с
  при длине кадров 1518 байт без потерь пакетов, пропущенных политикой.
  Проверка выполняется утилитой \texttt{iperf3}. На виртуальных стендах
  требование не проверяется. \der
\end{reqs}

\subsection{Требования к надёжности}

\begin{reqs}
\item[НДЖ-01] При аварийном завершении любой подсистемы ПО МЭ должно
  регистрировать событие и автоматически перезапускать подсистему.
\item[НДЖ-02] При неработоспособности подсистемы фильтрации передача трафика
  между интерфейсами не должна выполняться (отказ в «закрытое» состояние).
\item[НДЖ-03] После перезагрузки или восстановления электропитания ПО МЭ должно
  автоматически запускаться и применять сохранённую конфигурацию (УПР-10);
  содержимое журналов должно сохраняться (РЕГ-04).
\item[НДЖ-04] Все вводимые администратором данные должны проверяться
  на корректность; некорректные команды и значения отклоняются с выдачей
  сообщения без изменения состояния ПО МЭ.
\end{reqs}

\subsection{Условия эксплуатации}

Климатические условия эксплуатации определяются требованиями к используемым
техническим средствам (персональная ЭВМ или сервер общего назначения).

Эксплуатацию ПО МЭ осуществляют администраторы безопасности, обладающие
знаниями стека протоколов TCP/IP, принципов межсетевого экранирования
и навыками работы в ОС Linux. Количество администраторов не ограничивается.
Специальное обслуживание ПО МЭ не требуется.

\subsection{Требования к составу и параметрам технических средств}

Для функционирования ПО МЭ необходимо вычислительное устройство со следующими
минимальными характеристиками:
\begin{itemize}
\item процессор с архитектурой x86-64, не менее 2 ядер;
\item оперативная память не менее 2~ГБ;
\item накопитель не менее 8~ГБ;
\item от 1 до 6 сетевых адаптеров Ethernet 1~Гбит/с, поддерживаемых ядром ОС Linux;
\item средства локального управления: монитор и клавиатура или последовательная консоль.
\end{itemize}

Для разработки и испытаний допускается использование виртуальных машин
(QEMU/KVM, VirtualBox и т.\,п.) с виртуальными сетевыми адаптерами.

Для удалённого управления необходим компьютер администратора с SSH-клиентом
и web-браузером, имеющий сетевой доступ к одному из интерфейсов изделия.

\subsection{Требования к информационной и программной совместимости}

ПО МЭ должно функционировать под управлением ОС Linux с ядром версии 5.10
и выше (например, Debian~12, Ubuntu~22.04 и новее). Язык программирования
и используемые библиотеки выбираются разработчиком на этапе разработки
архитектуры; сторонние компоненты должны распространяться под свободными
лицензиями и устанавливаться из репозиториев дистрибутива.

ПО МЭ должно обрабатывать пакеты в соответствии со спецификациями протоколов
Ethernet~II, IPv4 (RFC~791), ICMPv4 (RFC~792), UDP (RFC~768) и TCP (RFC~9293).

Web-интерфейс должен корректно работать в актуальных версиях браузеров
Mozilla Firefox и Chromium (Google Chrome).

\subsection{Требования к маркировке и упаковке}

Требования не предъявляются. ПО МЭ поставляется в виде архива или пакета
дистрибутива, содержащего исполняемые файлы, конфигурацию по умолчанию,
сценарий установки и файл с указанием номера версии.

\subsection{Требования к транспортированию и хранению}

Требования не предъявляются. Хранение и передача ПО МЭ осуществляются
в электронном виде.

\subsection{Специальные требования}

Специальные требования не предъявляются.
\newpage
% ===================== СТАДИИ И ЭТАПЫ РАЗРАБОТКИ ======================
\section{Стадии и этапы разработки}

Разработка выполняется в стадии, приведённые в таблице~\ref{tab:stages},
в сроки, установленные графиком выполнения курсовой работы.

{\small
\begin{longtable}{|C{0.7cm}|L{3.5cm}|L{6.0cm}|L{3.9cm}|}
\caption{Стадии и этапы разработки}\label{tab:stages}\\
\hline
\textbf{№} & \textbf{Стадия (ГОСТ 19.102-77)} & \textbf{Содержание работ} & \textbf{Результат}\\
\hline\endfirsthead
\hline
\textbf{№} & \textbf{Стадия} & \textbf{Содержание работ} & \textbf{Результат}\\
\hline\endhead
1 & Техническое задание & Анализ требований задания, разработка спецификации требований & Техническое задание\\ \hline
2 & Эскизный проект & Разработка архитектуры ПО МЭ, разбиение на подсистемы и модули, выбор средств реализации & Описание архитектуры\\ \hline
3 & Рабочий проект & Реализация и отладка ПО МЭ & Текст программы, исполняемые файлы, руководство администратора\\ \hline
4 & Рабочий проект & Разработка программы и методики испытаний & Программа и методика испытаний\\ \hline
5 & Внедрение & Функциональное тестирование по программе и методике испытаний, защита работы & Протокол испытаний, пояснительная записка\\ \hline
\end{longtable}
}

% ===================== ПОРЯДОК КОНТРОЛЯ И ПРИЁМКИ =====================
\section{Порядок контроля и приёмки}

Контроль ПО МЭ осуществляется путём функциональных испытаний в соответствии
с программой и методикой испытаний, разработанной по ГОСТ~19.301-79.
Программа и методика испытаний должна содержать проверки, покрывающие все
требования подраздела~\ref{sec:func} настоящего документа (трассировка требований
задания~— в приложении~В).

Испытания проводятся на стенде, включающем изделие с ПО МЭ (не менее трёх
интерфейсов), не менее двух узлов в разных сегментах сети, подключённых
к разным интерфейсам изделия, и компьютер администратора. Для формирования
и анализа трафика допускается применение утилит \texttt{ping}, \texttt{nmap},
\texttt{hping3}, \texttt{scapy}, \texttt{tcpdump}, Wireshark, \texttt{curl}.

ПО МЭ считается выдержавшим испытания, если все проверки программы и методики
испытаний завершены успешно. Результаты оформляются протоколом испытаний.
Приёмка работы осуществляется руководителем курсовой работы.

% ============================ ПРИЛОЖЕНИЕ А ============================
\appsection{А}{Коды событий журнала событий}

{\small
\begin{longtable}{|C{1.2cm}|L{9.6cm}|C{3.5cm}|}
\caption{Коды событий}\label{tab:events}\\
\hline
\textbf{Код} & \textbf{Описание события} & \textbf{Требование}\\
\hline\endfirsthead
\hline
\textbf{Код} & \textbf{Описание события} & \textbf{Требование}\\
\hline\endhead
\multicolumn{3}{|l|}{\textit{1xx — аутентификация и сеансы администраторов}}\\ \hline
100 & Успешный вход администратора & РЕГ-01, ДСТ-01\\ \hline
101 & Неуспешная попытка входа & РЕГ-01, ДСТ-01\\ \hline
102 & Выход администратора & РЕГ-01\\ \hline
103 & Завершение сеанса по тайм-ауту неактивности & ДСТ-06\\ \hline
104 & Принудительный сброс сеанса администратора & УПР-12\\ \hline
105 & Понижение прав сеанса до «только чтение» & ДСТ-05\\ \hline
106 & Отказ в выполнении функции (недостаточно прав) & ДСТ-03\\ \hline
\multicolumn{3}{|l|}{\textit{2xx — включение и выключение МЭ и подсистем}}\\ \hline
200 & Запуск ПО МЭ & РЕГ-01\\ \hline
201 & Выключение изделия & РЕГ-01, УПР-14\\ \hline
202 & Перезагрузка изделия & РЕГ-01, УПР-14\\ \hline
210 & Запуск подсистемы & РЕГ-01\\ \hline
211 & Остановка подсистемы & РЕГ-01\\ \hline
212 & Перезапуск подсистемы & РЕГ-01, УПР-14\\ \hline
213 & Аварийное завершение подсистемы & РЕГ-01, НДЖ-01\\ \hline
\multicolumn{3}{|l|}{\textit{3xx — изменение политики доступа и правил NAT}}\\ \hline
300 & Добавление правила & РЕГ-01, УПР-03\\ \hline
301 & Изменение правила & РЕГ-01, УПР-03\\ \hline
302 & Удаление правила & РЕГ-01, УПР-03\\ \hline
303 & Перенумерация правил & РЕГ-01, УПР-03\\ \hline
304 & Активация / деактивация правила & РЕГ-01, УПР-03\\ \hline
305 & Добавление правила DNAT / SNAT & РЕГ-01, УПР-03, АПИ-09 \\ \hline
306 & Изменение правила DNAT / SNAT & РЕГ-01, УПР-03, АПИ-09 \\ \hline
307 & Удаление правила DNAT / SNAT & РЕГ-01, УПР-03, АПИ-09 \\ \hline
308 & Перенумерация правил DNAT / SNAT & РЕГ-01, УПР-03, АПИ-10 \\ \hline
309 & Активация / деактивация правила DNAT / SNAT & РЕГ-01, УПР-03 \\ \hline
\multicolumn{3}{|l|}{\textit{4xx — изменение конфигурационных параметров}}\\ \hline
400 & Включение / выключение интерфейса & УПР-13\\ \hline
401 & Изменение IP-адреса интерфейса & УПР-13\\ \hline
402 & Добавление системного маршрута & УПР-14\\ \hline
403 & Удаление системного маршрута & УПР-14\\ \hline
404 & Изменение даты и времени & УПР-14\\ \hline
405 & Включение / отключение механизма контроля состояний & УПР-08\\ \hline
406 & Изменение параметров контроля состояний & СОС-04\\ \hline
407 & Удаление записей таблицы состояний & УПР-08\\ \hline
408 & Сохранение текущей конфигурации & УПР-09\\ \hline
409 & Применение сохранённой конфигурации & УПР-09\\ \hline
410 & Применение конфигурации по умолчанию & УПР-09\\ \hline
411 & Очистка журнала пакетов & УПР-07, РЕГ-05\\ \hline
412 & Добавление учётной записи администратора & УПР-12\\ \hline
413 & Изменение учётной записи администратора & УПР-12\\ \hline
414 & Удаление учётной записи администратора & УПР-12\\ \hline
415 & Отклонена попытка удаления или понижения прав предустановленного администратора & ДСТ-04\\ \hline
416 & Изменение параметров регистрации, сеансов или контроля целостности & РЕГ-06, ДСТ-06, ЦЛС-03\\ \hline
\multicolumn{3}{|l|}{\textit{5xx — контроль целостности}}\\ \hline
500 & Проверка целостности завершена успешно & ЦЛС-03\\ \hline
501 & Выявлено нарушение целостности (с перечнем файлов) & ЦЛС-04\\ \hline
502 & Переход в режим аварийной поддержки, блокировка трафика & ЦЛС-04\\ \hline
503 & Восстановление файла из резервной копии & ЦЛС-05\\ \hline
504 & Выход из режима аварийной поддержки, возобновление обработки трафика & ЦЛС-06\\ \hline
505 & Обновление эталонных значений контрольных сумм & ЦЛС-02\\ \hline
\multicolumn{3}{|l|}{\textit{6xx — служебные события}}\\ \hline
600 & Достигнут максимальный размер журнала, удалены старые записи & РЕГ-06\\ \hline
\end{longtable}
}

% ============================ ПРИЛОЖЕНИЕ Б ============================
\appsection{Б}{Функции управления и уровни прав доступа}

Все функции таблицы~\ref{tab:functions} должны быть доступны как в командном,
так и в web-интерфейсе (УПР-02). Знак «+» означает, что функция доступна
администратору с данным уровнем прав.

{\small
\begin{longtable}{|C{0.9cm}|L{8.4cm}|C{1.9cm}|C{1.0cm}|C{1.0cm}|}
\caption{Функции управления}\label{tab:functions}\\
\hline
\textbf{№} & \textbf{Функция} & \textbf{Пункт задания} & \textbf{ЧЗ} & \textbf{ТЧ}\\
\hline\endfirsthead
\hline
\textbf{№} & \textbf{Функция} & \textbf{Пункт задания} & \textbf{ЧЗ} & \textbf{ТЧ}\\
\hline\endhead
\multicolumn{5}{|l|}{\textit{Политика доступа}}\\ \hline
1 & Просмотр правил фильтрации, DNAT и SNAT & 5.3 & + & +\\ \hline
2 & Добавление, изменение, удаление, перенумерация правил фильтрации, DNAT и SNAT, изменение активности & 5.3 & + & \\ \hline
\multicolumn{5}{|l|}{\textit{Журналы}}\\ \hline
3 & Просмотр журнала событий с выборкой по коду и интервалу времени & 5.4 & + & +\\ \hline
4 & Просмотр журнала пакетов с выборкой по критериям УПР-05 & 5.5 & + & +\\ \hline
5 & Очистка журнала пакетов & 5.6 & + & \\ \hline
\multicolumn{5}{|l|}{\textit{Контроль состояний}}\\ \hline
6 & Просмотр настроек механизма контроля состояний & 5.7 & + & +\\ \hline
7 & Включение / отключение механизма, изменение параметров & 5.7 & + & \\ \hline
8 & Просмотр таблицы состояний & 5.7 & + & +\\ \hline
9 & Удаление записей из таблицы состояний & 5.7 & + & \\ \hline
\multicolumn{5}{|l|}{\textit{Конфигурации}}\\ \hline
10 & Сохранение текущей конфигурации & 5.8 & + & \\ \hline
11 & Вывод текущей / сохранённой конфигурации & 5.8 & + & +\\ \hline
12 & Применение сохранённой конфигурации & 5.8 & + & \\ \hline
13 & Применение конфигурации по умолчанию & 5.8 & + & \\ \hline
\multicolumn{5}{|l|}{\textit{Учётные записи администраторов}}\\ \hline
14 & Добавление, удаление, редактирование учётной записи & 5.9 & + & \\ \hline
15 & Смена собственного пароля & 5.9 & + & +\\ \hline
16 & Просмотр базы учётных записей & 5.9 & + & +\\ \hline
17 & Просмотр активных сеансов & 5.9 & + & +\\ \hline
18 & Сброс активных сеансов & 5.9 & + & \\ \hline
\multicolumn{5}{|l|}{\textit{Интерфейсы}}\\ \hline
19 & Просмотр состояния интерфейсов & 5.10 & + & +\\ \hline
20 & Включение / выключение интерфейса & 5.10 & + & \\ \hline
21 & Настройка IP-адреса интерфейса & 5.10 & + & \\ \hline
\multicolumn{5}{|l|}{\textit{Системные настройки}}\\ \hline
22 & Просмотр даты и времени & 5.11 & + & +\\ \hline
23 & Настройка даты и времени & 5.11 & + & \\ \hline
24 & Просмотр системных маршрутов & 5.11 & + & +\\ \hline
25 & Добавление / удаление системных маршрутов & 5.11 & + & \\ \hline
26 & Выключение изделия & 5.11 & + & \\ \hline
27 & Перезагрузка изделия & 5.11 & + & \\ \hline
28 & Перезагрузка подсистемы ПО МЭ & 5.11 & + & \\ \hline
29 & Проверка целостности контролируемых файлов по запросу & 5.11, 6.1 & + & +\\ \hline
30 & Изменение параметров контроля целостности, регистрации, сеансов & 6.1 & + & \\ \hline
31 & Просмотр системной информации & 5.11 & + & +\\ \hline
\end{longtable}
}

% ============================ ПРИЛОЖЕНИЕ В ============================
\appsection{В}{Трассировка требований задания}

{\small
\begin{longtable}{|C{1.8cm}|L{7.2cm}|L{5.2cm}|}
\caption{Соответствие пунктов задания требованиям ТЗ}\label{tab:trace}\\
\hline
\textbf{Пункт задания} & \textbf{Содержание} & \textbf{Требования ТЗ}\\
\hline\endfirsthead
\hline
\textbf{Пункт задания} & \textbf{Содержание} & \textbf{Требования ТЗ}\\
\hline\endhead
1.1 & Функционирование на ОС Linux & ОБЩ-01\\ \hline
1.2 & От 1 до 6 портов Ethernet 1~Гбит/с & ОБЩ-02, ОБЩ-03\\ \hline
2.1 & Фильтрация всего трафика по политике доступа & ФЛТ-01, ФЛТ-02, ФЛТ-03, ФЛТ-12\\ \hline
2.2 & Критерии фильтрации & ФЛТ-04, ФЛТ-05, ФЛТ-06\\ \hline
2.3 & Атрибуты правил & ФЛТ-07, ФЛТ-08, ФЛТ-09, СОС-01 -- СОС-04\\ \hline
2.4 & Source NAT и Destination NAT & ФЛТ-10, ФЛТ-11, НАТ-01 -- НАТ-16, СТР-01 -- СТР-08 \\ \hline
3.1 & Регистрируемые события & РЕГ-01, РЕГ-04, прил.~А\\ \hline
3.2 & Состав записи журнала событий & РЕГ-02\\ \hline
3.3 & Состав записи журнала пакетов & РЕГ-03, РЕГ-04\\ \hline
4.1 & Аутентификация по идентификатору и паролю & ДСТ-01, ДСТ-02, ДСТ-06, ДСТ-07\\ \hline
4.2 & Уровни прав ЧЗ и ТЧ & ДСТ-03, прил.~Б\\ \hline
4.3 & Предустановленный администратор & ДСТ-04\\ \hline
4.4 & Понижение прав второго администратора ЧЗ & ДСТ-05\\ \hline
5.1 & Локальное и удалённое управление & УПР-01\\ \hline
5.2 & Командный и web-интерфейс & УПР-02\\ \hline
5.3 & Просмотр и редактирование правил, API управления политикой и NAT & УПР-03, АПИ-01 -- АПИ-12\\ \hline
5.4 & Просмотр и выборка журнала событий & УПР-04, УПР-06\\ \hline
5.5 & Просмотр и выборка журнала пакетов & УПР-05, УПР-06\\ \hline
5.6 & Очистка журнала пакетов & УПР-07, РЕГ-05\\ \hline
5.7 & Управление контролем состояний & УПР-08, СОС-03, СОС-05\\ \hline
5.8 & Управление конфигурациями & УПР-09, УПР-10, УПР-11\\ \hline
5.9 & Управление учётными записями & УПР-12\\ \hline
5.10 & Настройка Ethernet-интерфейсов & УПР-13\\ \hline
5.11 & Системные настройки и режимы & УПР-14, УПР-15, ОБЩ-04\\ \hline
6.1 & Периодический контроль целостности и по запросу & ЦЛС-01, ЦЛС-02, ЦЛС-03\\ \hline
6.2 & Действия при нарушении целостности & ЦЛС-04 -- ЦЛС-07\\ \hline
2.1 -- 2.4, 1.2 & Реализация приёма, обработки и передачи Ethernet-кадров и IP-пакетов в ОС Linux & ОБР-01 -- ОБР-13\\ \hline
\end{longtable}
}

\end{document}